\section{Introduction}
\label{sec:introduction}
Exact mixed-integer linear programming (MILP) solvers can emit proofs in the VIPR format, whose inferences are checked in rational arithmetic independently of the solver \citep{cheung2017-vipr,hojny2025-the-scip-optimization-suite-10}. Such a proof certifies a bound for its own input model. For a convex mixed-integer nonlinear program (MINLP), outer approximation (OA) supplies the classical link to a MILP: affine underestimators of the convex functions define a MILP master whose lower bounds transfer to the original problem \citep{duran1986-an-outer-approximation-algorithm-for,kronqvist2019-a-review-and-comparison-of}. Certifying an MINLP bound through this link requires two further checks. Every nonlinear cut must be valid for the stated nonlinear model, and the proved MILP must be exactly the master justified by those cuts and the original constraints. Without these checks, exact discrete arithmetic can certify the wrong relaxation. For example, rounding a tangent's slope without correcting its intercept can exclude a feasible point (Example~\ref{ex:rounding}), and a proof for a different master cannot repair that error.

For minimization, we seek a rational number $\beta$ and evidence establishing
\begin{equation}
  \beta\leq f(x)\qquad\text{for every feasible mixed-integer point }x
  \label{eq:intro-bound}
\end{equation}
for a precisely interpreted nonlinear model. The evidence should be checkable without repeating the numerical search that produced it.

We develop and implement a certificate interface that performs both checks. One exact interpretation of the loaded expression tree governs domain and curvature checks, derivative and interval evaluation, and master construction. Rigorous enclosures at a supplied support point justify each rational affine cut. The checker regenerates the rational master, matches it to the input of a VIPR proof, and replays every discrete inference in rational arithmetic. Numerical search supplies candidate points, slopes, and proofs; acceptance depends only on their checked properties, not on the accuracy of numerical subproblem solutions.

\subsection{Contribution and scope}
\label{sec:contribution}
The contribution is the integration of established convex-support and discrete-proof principles into a checked nonlinear-to-linear certificate interface, together with evidence about its reliability.
\begin{enumerate}
\item \emph{Soundness.} Accepted evidence satisfying the stated contracts implies \eqref{eq:intro-bound} (Theorem~\ref{thm:transfer}). The proof combines three parts: rational bound propagation preserves feasible points (Lemma~\ref{lem:propagation}); a rational cut test covers bounded, one-sided, and free coordinates (Theorem~\ref{thm:safecut}, Corollary~\ref{cor:enclosure}); and an invariant for the admitted VIPR inferences gives an unconditional master bound (Proposition~\ref{prop:vipr-invariant}, Theorem~\ref{thm:master-bound}). The argument maps every feasible point into the master with the same objective value. It needs no minimizer, Slater point, or finite optimum. An independently checked feasible point turns the bound into a certified optimality gap (Corollary~\ref{cor:primal}).
\item \emph{Implementation.} The checker recognizes curvature through sufficient composition rules and exact tests for quadratics, monomials, and one-variable linear fractions (Section~\ref{sec:curvature-rules}). It checks cuts with outward interval arithmetic and exact endpoint extraction, binds the proof to the regenerated master under a checked variable bijection, and accepts complete proofs in a restricted subset of VIPR (Sections~\ref{sec:cut-implementation}--\ref{sec:proof-implementation}). It rejects two unsound behaviors reproduced in an inspected version of the external checker \texttt{viprchk} (Section~\ref{sec:checker-repairs}). A suite of 161 targeted tests exercises the executable contracts and regression cases.
\item \emph{Formal verification.} Lean proofs cover 49 mathematical obligations: propagation, curvature and cut rules, typed discrete checking, master identity, bound transfer, and primal completion. They include executable reference checkers for typed discrete proofs and rational quadratic matrices (Section~\ref{sec:formalization}). The nonlinear theorems take genuine derivative and enclosure data as hypotheses. The Python expression, interval, master-matching, and proof-checking software remains in the trusted base: it is not shown to refine the Lean checkers, and no benchmark certificate is checked in Lean.
\item \emph{Experiments and audits.} On 289 convex MINLPLib models \citep{minlplib}, a frozen uniform protocol with 30-second search budgets yields artifacts that pass separate replay for 203 models. Of the rest, 19 are rejected and 67 fail curvature or domain admission before a master is built (Section~\ref{sec:replay-results}). Of the 203 bounds, 80 lie at or below an unverified library reference value, with normalized difference at most $10^{-2}$. Exact audits (Section~\ref{sec:failure-mechanisms}) find three newly generated proofs that the external checker accepted although they contain invalid linear-combination or integer-disjunction steps, and show that two saved solver points are infeasible. A checked feasible point completes an exact optimality proof for \texttt{clay0204m}. Accepted uniform-run proofs total about 14.2\,GB (Section~\ref{sec:experimental-costs}).
\item \emph{Reusable artifacts.} Portable checker code, frozen inputs and outcomes, complete proof bundles, exact failure extracts, and original-model witness audits support reproduction (Appendix~\ref{app:reproduction}). A descriptive catalogue collects checked bounds for 222 distinct loaded models across all recorded protocols; it is a collection, not a uniform-run success rate.
\end{enumerate}
We claim neither a new general support-minimization principle nor a certificate-size or solver-performance advantage over prior work.

The implementation covers supported smooth convex expressions and exactly recognized quadratic and structured functions. General nonconvex MINLP, nonlinear equalities, and unsupported expression domains lie outside its contract. It certifies the exact interpretation of the loaded expression model, whose relation to the source modeling file is specified in Section~\ref{sec:model}. Checked feasible points exist for \texttt{clay0204m} and a small quadratic example; proximity to a library reference value supplies no such point.

\subsection{Relationship to prior work}
\label{sec:prior-work}
\paragraph{Mixed-integer convex optimality certificates.}
Certificates for convex MINLP are established. \citet{baes2016-duality-for-mixed-integer-convex} derive finite optimality certificates and mixed-integer duality using subgradients and mixed-integer-free polyhedra. Under their respective assumptions, their Theorems 1 and 3 bound the number of certificate points by $2^{|J|}$, where $|J|$ is the number of integer variables; this bounds the number of points, not the bit length of their real-valued data. \citet{basu2017-optimality-certificates-for-convex-minimization} relate finite convex optimality certificates to Helly numbers for more general feasible sets.

The closest computational work is \citet{halbig2024-computing-optimality-certificates-for-convex}. Their generalized certificates consist of suspension points (points that anchor the certificate hyperplanes) and hyperplanes. They give construction and reduction algorithms, structural bounds on certificate size, and experiments on MINLPLib and MIPLIB. Their verification removes integrality and solves a continuous convex problem to test each certificate plane (Algorithm~3); a linear mixed-integer problem tests integer-freeness of the certificate polyhedron.

Our interface instead checks each cut by recomputing rigorous function and derivative enclosures at a supplied support point, then replays saved rational inferences for the discrete bound. The result is an OA certificate with a different checking procedure from the geometric certificates above. The certified bound may come from an interrupted master solve and need not equal a known optimum. Theorem~\ref{thm:transfer} is a conditional validity result, not a completeness theorem guaranteeing such evidence for every optimum.

\paragraph{Outer approximation and numerical safeguards.}
\citet{coey2020-outer-approximation-with-conic-certificates} derive OA cuts from conic dual certificates and analyze scaling under MILP feasibility tolerances. We instead check affine underestimation for explicitly interpreted nonlinear expressions and export an exact rational master proof; the two approaches differ in their nonlinear representations and checking obligations.

Directed rounding and rational repair are established safeguards for approximate numerical information. \citet{neumaier2004-safe-bounds-in-linear-and} develop safe bounds and cuts for linear and mixed-integer linear programming, \citet{eifler2024-safe-and-verified-gomory-mixed} combine safe Gomory mixed-integer cuts with rational proof verification, and \citet{borst2024-certified-constraint-propagation-and-dual} treat certified propagation and dual proof analysis in an exact MIP solver.

For nonlinear problems, \citet{jansson2004-rigorous-convex} develops rigorous lower bounds for convex optimization, and \citet{jansson2009-verified-convex} surveys verified convex and conic computations and their use in reliable global and combinatorial optimization. \citet{messine2019-reliable-relaxation} minimize supporting affine functions over finite boxes to bound convex relaxations; our finite-box correction applies this calculation to the residual left after subtracting the proposed affine row. \citet{elloumi2025-qibex} describe QIBEX-R, which combines interval branch-and-bound for quadratic mixed-integer problems with curvature correction and rigorous relaxation bounds. These methods are precedents for our nonlinear layer. Our contribution joins that layer to explicit model and master binding and to replay of a saved discrete proof.

\paragraph{MILP proofs and checker correctness.}
\citet{cheung2017-vipr} introduce the VIPR format, whose inferences can be checked independently of the MILP solver that produced them. Exact MILP solving and VIPR proof production are available in SCIP \citep{hojny2025-the-scip-optimization-suite-10}, and recent work constructs VIPR proofs from a black-box ILP solver by rational reconstruction \citep{szeider2026-blackbox-vipr}. These methods supply discrete evidence; they do not establish the validity of nonlinear cuts supplied as input.

The semantics of a proof format and the correctness of a particular checker are separate issues. \citet{wood2026-vipr-smt} prove a logical encoding of VIPR validity correct in Why3 and implement an SMT-based checker; their theorem covers the logical transformation, not the C++ translation program. They also identify trust-sensitive aspects of the existing checker. The CakeML project contains a VIPR formalization, a pure checker, and a verified program with theorems about its file and standard-input semantics \citep{cakeml2026-vipr}. Our Python replay kernel is a separate, unmechanized implementation for a restricted subset of VIPR. Its acceptance contract additionally binds the proof to the input master and rejects unsupported or malformed cases.

\paragraph{Formal nonlinear optimization and verified reductions.}
\citet{solovyev2013-formal-verification-of-nonlinear-inequalities} verify nonlinear inequalities with Taylor interval approximations in HOL Light. \citet{narkawicz2014-a-formally-verified-generic-branching} verify a generic global-optimization branching algorithm in PVS. \citet{magron2015-formal} combine semialgebraic approximations and sums-of-squares witnesses with Coq to prove nonlinear lower bounds. These are precedents for proof-producing numerical computation, with different representations of the nonlinear and discrete reasoning. CvxLean verifies optimization reductions in Lean \citep{bentkamp2023-verified-reductions-for-optimization}; its proofs show that a transformed problem represents the original one, while the external numerical solve is a separate obligation. The same distinction governs our formalization (Section~\ref{sec:formalization}).

\paragraph{Organization.}
Section~\ref{sec:model} fixes the model and the certificate contract. Section~\ref{sec:soundness} proves soundness. Section~\ref{sec:implementation} describes the executable checker and Section~\ref{sec:formalization} the Lean coverage. Section~\ref{sec:experiments} reports the experiments and audits, and Section~\ref{sec:discussion} concludes. Appendix~\ref{app:reproduction} describes the accompanying materials, and Appendix~\ref{app:experimental-accounting} records the experimental protocols and full accounting.
